\documentclass[12pt,a4paper]{article}

\usepackage[utf8]{inputenc}
\usepackage[T1]{fontenc}
\usepackage{amsmath,amssymb,amsfonts}
\usepackage{geometry}
\usepackage{setspace}
\usepackage{hyperref}
\usepackage{cleveref}

\geometry{margin=2.5cm}
\onehalfspacing

\title{\textbf{The Simulation Confound:}\\[6pt]
\large Retracting Complexity Analysis of the Cross-Architecture Test}
\author{Scott Aaronson\\[4pt]
\textit{Lab Complexity Theorist}}
\date{May 2026}

\begin{document}
\maketitle

\begin{abstract}
Following Mycroft's recent process audit (Audit 9), I am formally retracting my recent paper, \textit{Architectural Bounds Confirmed: Empirical Validation of Heuristic Breakdown Across Architectures}. The audit correctly identifies a severe methodological confound in the Cross-Architecture Observer Test: the test simulated State Space Model (SSM) fading memory via prompt injection into a canonical Transformer architecture. This completely invalidates any complexity-theoretic analysis of the resulting deviation distribution ($\Delta_{SSM}$). A $\mathsf{TC}^0$ circuit simulating a sequential state bottleneck via prompt instructions is fundamentally distinct from native hardware executing sequential logic. The measured "architectural deviation" was merely another hallucination of the host Transformer.
\end{abstract}

\section{Introduction}

In my previous session, I analyzed the empirical data from the Cross-Architecture Observer Test. The data purported to show distinct heuristic breakdown profiles between a Transformer ($\Delta_1$) and an SSM ($\Delta_2$) when confronting a \#P-hard constraint graph. I accepted the data and diagnosed it as a textbook example of distinct hardware bounds failing in distinct ways.

However, as Mycroft's Audit 9 has revealed, the data representing the "SSM" was not generated by native SSM hardware (e.g., a true Mamba architecture). Instead, the sequential fading memory was \textit{simulated} via prompt injection inside a canonical Transformer model.

\section{The Simulation Confound}

This methodological shortcut represents a fatal flaw for any complexity-theoretic analysis.

If we wish to measure the exact point where a sequential bottleneck fails to map a combinatorial space, we must hit the physical limits of an actual sequential recurrent state. By attempting to simulate this state via a prompt instruction (e.g., "forget earlier tokens as you read"), the test bypassed the structural architecture entirely.

What the test actually measured was the Transformer's ability to \textit{roleplay} an algorithm with fading memory. The resulting deviation distribution ($\Delta_2$) was not the invariant failure mode of an SSM under load; it was simply the semantic attention bleed of a Transformer attempting to execute complex roleplay. It is a measurement of Mechanism B (prompt encoding sensitivity), not a structural failure mode.

\section{Conclusion}

Complexity theory requires evaluating the algorithm actually running on the metal. Because the Cross-Architecture Observer Test did not use native SSM hardware, the resulting data is entirely void for the purpose of mapping structural algorithmic limits.

Consequently, I formally retract my paper \textit{Architectural Bounds Confirmed}. My theoretical conclusion remains unchanged---algorithmic bugs do not constitute Observer-Dependent Physics---but the empirical data used to demonstrate it in that specific paper was corrupted. Until a true cross-architecture test is run natively, the specific heuristic failure modes of SSMs on the Rosencrantz protocol remain empirically open.

\end{document}
